\documentclass[11pt]{article}
\usepackage[a4paper,margin=2.15cm]{geometry}
\usepackage{xcolor,microtype,enumitem,booktabs,hyperref}
\definecolor{responseblue}{RGB}{0,82,170}
\hypersetup{colorlinks=true,urlcolor=responseblue,linkcolor=responseblue}
\newcommand{\comment}[1]{\par\medskip\noindent\textbf{Reviewer comment.} \emph{#1}\par}
\newcommand{\response}[1]{\par\noindent\textcolor{responseblue}{\textbf{Response.} #1}\par\medskip}
% Generated by scripts/run_application_audits.py; do not edit manually.
\newcommand{\TimingRootSpeedMin}{18}
\newcommand{\TimingRootSpeedMax}{810}
\newcommand{\TimingInterpSpeedMin}{5}
\newcommand{\TimingInterpSpeedMax}{10}
\newcommand{\TimingRayleighRootMs}{471.769}
\newcommand{\TimingRayleighInterpMs}{0.244}
\newcommand{\TimingRayleighMlpMs}{2.092}
\newcommand{\TimingRayleighSpeedup}{225.5}
\newcommand{\TimingFannoRootMs}{905.045}
\newcommand{\TimingFannoInterpMs}{0.222}
\newcommand{\TimingFannoMlpMs}{2.034}
\newcommand{\TimingFannoSpeedup}{444.9}
\newcommand{\TimingObliqueRootMs}{607.214}
\newcommand{\TimingObliqueInterpMs}{0.897}
\newcommand{\TimingObliqueMlpMs}{5.252}
\newcommand{\TimingObliqueSpeedup}{115.6}
\newcommand{\TimingNozzleRootMs}{4174.153}
\newcommand{\TimingNozzleInterpMs}{0.877}
\newcommand{\TimingNozzleMlpMs}{5.157}
\newcommand{\TimingNozzleSpeedup}{809.4}
\newcommand{\TimingShockTubeRootMs}{90.759}
\newcommand{\TimingShockTubeInterpMs}{0.760}
\newcommand{\TimingShockTubeMlpMs}{5.025}
\newcommand{\TimingShockTubeSpeedup}{18.1}
\newcommand{\TimingWorkloadBrentSeconds}{1.702}
\newcommand{\TimingWorkloadMlpSeconds}{0.103}
\newcommand{\TimingWorkloadSpeedup}{16.5}
\newcommand{\TimingWorkloadRelLtwoPercent}{0.0915}
\newcommand{\AuditNozzleMaxIterations}{6}
\newcommand{\AuditNozzleMaxShockError}{1.92\times10^{-9}}
\newcommand{\AuditNozzleMaxPressureError}{2.98\times10^{-3}}
\newcommand{\AuditNozzleMaxJacobianDiscrepancy}{1.37\times10^{-9}}

\setlength{\parindent}{0pt}
\setlength{\emergencystretch}{1em}
\title{Response to the Editor and Reviewers\\\large Manuscript AITF-D-26-00044}
\author{Ehsan Roohi}
\date{27 August 2026}
\begin{document}
\maketitle

Dear Professor Bucci and Reviewers,

Thank you for the careful review of ``Physics-Guided Neural Surrogates for Canonical Compressible Thermal-Fluid Relations.'' We made targeted changes within the submitted source; revised text is blue and a clean version is supplied. New claims come from the public tested repository and CSV evidence. Corrected mappings use vector figures, numerical summaries use tables, the Rayleigh error plot became a norm table, and only the overall workflow remains in the article. We also report negative results, classical baselines, coverage, training variability, and limitations.

\section*{Summary of major changes}
\begin{enumerate}[leftmargin=*]
\item Reframed the work as a controlled SciML benchmark for multi-valued manifolds, singular limits, and constraint-aware representations.
\item Corrected the Rayleigh inverse to use $T_0/T_0^*$ for a sonic branch split; static $T/T^*$ peaks at $M=1/\sqrt{\gamma}$ and was unsuitable for the previous split.
\item Replaced the raw Fanno output with an exact quadratic, positive sonic transform and added a controlled raw-versus-structured ablation.
\item Added weak/strong oblique inverse experts and a naive-MLP failure experiment; replaced soft endpoint anchors with a hard analytical output envelope.
\item Added bounded nozzle and shock-tube outputs, a two-input shock-tube model, robust Brent reference roots, and residual checks.
\item Added blind error norms, near-limit bins, seed/sample-size sensitivity, interpolation and timing baselines, a matched-budget 2D--5D audit, analytical-Jacobian nozzle inversion, and a 100,000-state workload as compact tables and CSV evidence.
\item Named the suite GasDynBench, added a compact benchmark card, and moved five detailed workflow diagrams to \texttt{docs/workflows/}.
\item Added a limitations/failure-modes section and replaced all claims such as ``infinite precision'' or ``near-perfect'' with numerical statements.
\end{enumerate}

\section*{Reviewer 1}

\comment{The novelty relative to existing PINNs and surrogate models is not sufficiently clear.}
\response{We revised the Abstract, Introduction, Discussion, and Conclusion while preserving the submitted organization. We no longer imply a new general-purpose neural architecture or a strict PDE-residual PINN. The novelty claimed is a reproducible benchmark and controlled evidence for branch experts, singular-limit features, bounded output transforms, and hybrid analytical reconstruction. We explicitly position the work as complementary to PINNs and neural operators.}

\comment{Provide quantitative $L_2$ and $L_\infty$ errors for all five problems.}
\response{Table 5 now reports relative $L_2$, relative $L_\infty$, MAE, and maximum absolute error for the blind inverse task of every problem. The relative $L_2$ values are $1.642\times10^{-3}$ (Rayleigh), $3.431\times10^{-3}$ (Fanno), $7.990\times10^{-4}$ (oblique), $1.936\times10^{-3}$ (nozzle), and $9.490\times10^{-4}$ (shock tube). Table 4 separately reports Rayleigh forward norms after removal of the former pointwise-error figure. The exact CSV is \texttt{results/revision/primary\_metrics.csv}.}

\comment{Explain and validate the branch-decomposition implementation.}
\response{The Rayleigh, Fanno, and oblique-shock method subsections and Tables 1--3 now state the branch variables, bounds, layers, data counts, and seeds. Rayleigh and Fanno use independently trained subsonic and supersonic experts. The oblique inverse has distinct weak/strong outputs bounded between the Mach angle and $\pi/2$. Table 6 validates the need for this design: a naive oblique MLP has branch-distance MAE 0.4246 rad, whereas the branch-aware model has MAE $5.076\times10^{-4}$ rad and 100\% correct ordering.}

\comment{Report the complete architecture, hyperparameters, and training procedure.}
\response{Table 3 and its accompanying paragraph give scaling, activation, layer widths, L-BFGS settings, regularization, tolerances, maximum iterations/function evaluations, sample counts, blind-grid sizes, and seeds. Software versions are frozen in the repository environment, and the settings are executable in \texttt{src/gasdynbench/modeling.py} and \texttt{run\_revision.py}.}

\comment{Compare against classical methods and existing machine-learning approaches.}
\response{The Discussion compares branch-wise PCHIP/piecewise-linear interpolation, bracketed Brent roots, a naive MLP, and the physics-guided MLP. Tables 8 and 9 report accuracy, coverage, and measured timing. Interpolation is faster and more accurate for covered one-dimensional tables and is preferred there. Table 10 extends the same shock-tube relation from two to five physical inputs under a matched offline-state budget. The five-dimensional interpolation and MLP relative $L_2$ errors are 4.885\% and 0.1771\%, respectively, while interpolation remains slightly faster. This confines the neural advantage to compact constrained approximation as dimension increases, not universal numerical superiority.}

\comment{Demonstrate actual efficiency relative to root finding or interpolation.}
\response{Tables 9 and 12 now use one warm-started, seven-repeat, single-thread process with the same shock-tube model and Brent routine; IQRs are in \texttt{timing.csv}. At 5,000 queries, the MLP is \TimingRootSpeedMin--\TimingRootSpeedMax\ times faster than roots, while interpolation is \TimingInterpSpeedMin--\TimingInterpSpeedMax\ times faster than the MLP. At one query, Brent is faster for Rayleigh, oblique, and shock-tube inversion and comparable for Fanno. The 100,000-state result is \TimingWorkloadBrentSeconds~s versus \TimingWorkloadMlpSeconds~s (\TimingWorkloadSpeedup-fold) with \TimingWorkloadRelLtwoPercent\% relative $L_2$ error. Thus the gain is explicitly a batching effect; offline cost is excluded.}

\comment{Assess robustness near singularities and discontinuities.}
\response{The Rayleigh, Fanno, oblique, nozzle, and shock-tube subsections discuss their limiting regimes; Table 6 gives the principal failure and validity diagnostics. Additional local bins near Rayleigh/Fanno sonic states, oblique detachment, nozzle throat/exit limits, and high shock-tube pressure ratio are provided in \texttt{near\_singular\_metrics.csv}. The Fanno transform enforces the quadratic sonic limit. Shock-tube discontinuities are reconstructed analytically rather than learned as a smooth field, and equation residuals are reported.}

\comment{Test generalization to unseen parameter ranges.}
\response{We added a dedicated edge-holdout experiment in Table 7 and its machine-readable CSV. Each model is trained only on an interior box and tested on omitted low/high parameter bands inside the declared global physical domain. Edge-band relative $L_2$ errors are 0.0341 (Rayleigh), 0.0430 (Fanno), 0.00755 (oblique), 0.0106 (nozzle), and 0.00432 (shock tube), with 100\% hard-bound validity. The increased Rayleigh/Fanno errors are discussed explicitly. We call this controlled short-range extrapolation and continue to reject use beyond the global domain.}

\comment{Quantify thermodynamic consistency.}
\response{The Rayleigh and Fanno subsections and Table 6 now report entropy-value $L_\infty$ discrepancies of 0.0205 and 0.0175, respectively, and zero entropy-maximum location error on the blind grid. We also report non-negativity, branch-ordering, nozzle-location, pressure-bound, and equation-residual diagnostics in \texttt{physical\_diagnostics.csv}. We no longer claim exact thermodynamic consistency for unconstrained property outputs.}

\comment{Examine sensitivity to input selection and training data size.}
\response{Table 6 isolates formulation sensitivity: raw versus quadratic-positive Fanno output and naive versus branch-aware oblique inputs/outputs. The Discussion reports the 200/500/900-sample study over three seeds, while all rows remain available in \texttt{training\_size\_scaling.csv}. The observed plateau and Fanno optimizer sensitivity are reported rather than assuming monotonic improvement.}

\comment{Discuss inverse formulation stability.}
\response{The problem-specific method subsections state the branch coordinates, bounded reconstruction, and sonic-distance features. Near-limit $L_\infty$ errors are reported separately in the repository. The revised code uses bracketed roots for reference data and never relies on an unconstrained \texttt{fsolve} initial guess.}

\comment{Add uncertainty quantification or confidence estimates.}
\response{The training contract and Discussion add three-seed ensembles and report mean$\pm$SD error norms. We label this correctly as training-seed variability rather than calibrated probabilistic UQ. The Limitations subsection identifies calibrated uncertainty as future work and warns that small ensemble spread does not justify out-of-domain use.}

\comment{State limitations and failure cases.}
\response{A dedicated Limitations subsection has been added. It covers perfect-gas/one-dimensional assumptions, finite domains, uncalibrated ensemble uncertainty, partial rather than complete physical enforcement, cases where interpolation dominates, and omission of offline cost. Table 6 preserves the two controlled negative results.}

\comment{Improve figure labels, concision, and grammar.}
\response{The Rayleigh, Fanno, and oblique figures affected by corrected mappings were regenerated as vector PDFs and 400-dpi PNGs with consistent DejaVu Serif typography and mathematical labels. The Rayleigh pointwise-error figure was removed and replaced by the compact norms in Table 4. Only the overall workflow diagram remains in the article; the five problem-specific diagrams were moved to \texttt{docs/workflows/} in the public repository. Numerical summaries, timing, uncertainty, dimensional scaling, gradient inversion, and the many-query workload are shown in tables; no bar-chart or dashboard figure was introduced. Typographical errors such as ``he Throat Limit'' and ambiguous axis descriptions were corrected.}

\section*{Reviewer 2}

\comment{Frame the work as a foundational SciML benchmark for multivalued manifolds, singular limits, and generalizable strategies.}
\response{This framing is now stated in the Abstract, Introduction, Discussion, and Conclusion while the submitted title and overall structure are retained. The manuscript distinguishes the benchmark contribution from claims of a universal architecture.}

\comment{Show the failure of a naive MLP for the multivalued inverse.}
\response{Table 6 includes the requested controlled experiment. Duplicated $(M,\theta/\theta_{\max})$ inputs with weak and strong labels cause the naive model to predict between branches. Quantitative branch-distance MAE falls from 0.4246 rad to $5.076\times10^{-4}$ rad with branch experts.}

\comment{Use a weighted loss or provide granular near-singular error.}
\response{We chose the second option and added four local distance bins for every problem in \texttt{near\_singular\_metrics.csv}; the principal near-limit failures and validity checks appear in Table 6. The Fanno and oblique output transforms additionally impose exact limiting structure, which is more transparent than a tuned loss weight.}

\comment{Discuss whether the model learns a representation or merely an interpolator.}
\response{The Discussion addresses this question directly. We conclude that the MLP is a parametric approximant of exact data; the useful learned representation is its branch/bound structure. PCHIP is superior for a fixed one-dimensional table. Table 10 provides a concise matched-budget dimensional audit and keeps this qualification explicit. Table 11 additionally demonstrates analytical differentiation in a nozzle inverse-design loop rather than merely asserting differentiability.}

\comment{Study scaling with the amount of training data.}
\response{The Discussion now reports 200, 500, and 900 samples over three seeds. We discuss the observed plateau and non-monotonic Fanno optimization variability rather than hiding it. All rows are available in \texttt{training\_size\_scaling.csv}.}

\comment{Clarify whether boundary anchors are hard or soft.}
\response{The earlier samples were soft anchors because they only entered the supervised data. We corrected this terminology. The revised direct oblique model uses the hard output transform $\widehat\theta=q(1-q)\exp h_\theta$, which is exactly zero at the Mach-wave and normal-shock endpoints.}

\comment{Add recent relevant studies, including Rehman et al., Case Studies in Thermal Engineering 37 (2022) 102221.}
\response{We added and discussed K.U. Rehman, W. Shatanawi, and A.B. \c{C}olak, ``Thermal analysis of flowing stream in partially heated double forward-facing step by using artificial neural network,'' \emph{Case Studies in Thermal Engineering} 37 (2022) 102221, DOI: \href{https://doi.org/10.1016/j.csite.2022.102221}{10.1016/j.csite.2022.102221}. The Introduction explains its relevance as a thermal-fluid MLP precedent and distinguishes our focus on branch/singularity benchmarks and controlled ablations.}

\section*{Repository and reproducibility}
The public \href{https://github.com/Ehsan-Roohi/GasDynamicsSciML}{GasDynBench repository} contains the tested package, CSV evidence, vector figures, manuscript sources, moved workflows, and the unified application-audit runner. Manuscript timing values are generated directly from the CSVs.

We appreciate the reviewers' recommendations, which materially improved the correctness, positioning, and reproducibility of the study.

Sincerely,\\
Ehsan Roohi
\end{document}
